\section{The Mermin relaxation-time dielectric function}\label{app:D.9}

This section derives the Mermin response of \cref{sec:4.5}~\cite{mermin1970lindhard}, in which collisions at the rate $\gamma_{\mathrm{rel}}$ relax the electrons toward a local equilibrium that has the instantaneous density of the electrons, and compares it with the substitution $\omega \to \omega + \ii\gamma_{\mathrm{rel}}$ in the independent-electron response.

\subsection*{The equation of motion}

Take independent electrons, $H_0 = \sum_{\vect{k}}\xi_{\vect{k}}c^\dagger_{\vect{k}}c_{\vect{k}}$ with the spin in the label, in equilibrium with occupations $n_{\vect{k}} = n_{\mathrm{F}}(\xi_{\vect{k}})$, and couple a potential to the density as in \cref{eq:4.1-kubo}, $H(t) = H_0 + n_e(\vect{q})\,\phi_{\mathrm{ext}}(t)$, with $\vect{q} \ne 0$. In the random-phase approximation of \cref{sec:4.5} the potential becomes the total one, external plus induced; here it is simply the potential the electrons feel. The response is carried by the coherences
\begin{equation*}
  \delta\varrho_{\vect{k}}(t) = \langle c^\dagger_{\vect{k}}c_{\vect{k}+\vect{q}}\rangle(t) ,
\end{equation*}
the matrix elements of the one-electron density matrix between $\vect{k} + \vect{q}$ and $\vect{k}$, which vanish in equilibrium. Their sum is the induced density, $\delta\nd(t) \equiv \Vol^{-1}\sum_{\vect{k}}\delta\varrho_{\vect{k}} = \Vol^{-1}\delta\langle n_e(\vect{q})^\dagger\rangle$, since $n_e(\vect{q})^\dagger = \sum_{\vect{k}}c^\dagger_{\vect{k}}c_{\vect{k}+\vect{q}}$ (\cref{eq:2.3-nq}). By \cref{eq:2.3-bilinear}, $[c^\dagger_{\vect{k}}c_{\vect{k}+\vect{q}}, H_0] = (\xi_{\vect{k}+\vect{q}} - \xi_{\vect{k}})\,c^\dagger_{\vect{k}}c_{\vect{k}+\vect{q}}$ and $[c^\dagger_{\vect{k}}c_{\vect{k}+\vect{q}}, n_e(\vect{q})] = c^\dagger_{\vect{k}}c_{\vect{k}} - c^\dagger_{\vect{k}+\vect{q}}c_{\vect{k}+\vect{q}}$. To first order in the potential the second commutator can be averaged in equilibrium, and $\ii\,\partial_t\langle A\rangle = \langle[A, H]\rangle$ gives
\begin{equation*}
  \ii\,\partial_t\,\delta\varrho_{\vect{k}} = -D_{\vect{k}}\,\delta\varrho_{\vect{k}} + (n_{\vect{k}} - n_{\vect{k}+\vect{q}})\,\phi_{\mathrm{ext}}(t), \qquad D_{\vect{k}} = \xi_{\vect{k}} - \xi_{\vect{k}+\vect{q}} .
\end{equation*}
For $\phi_{\mathrm{ext}} \propto \ee^{-\ii\omega t}$ the solution is $\delta\varrho_{\vect{k}} = (n_{\vect{k}} - n_{\vect{k}+\vect{q}})\,\phi_{\mathrm{ext}}/(\omega + \ii0^+ + D_{\vect{k}})$, and summing over $\vect{k}$ gives back $\delta\nd = R^0(\vect{q},\omega)\,\phi_{\mathrm{ext}}$ with the bubble of \cref{eq:4.7-bubble}. The equation of motion is the one-electron Liouville equation, and collisions add a relaxation term to it.

\subsection*{Relaxation toward local equilibrium}

In the relaxation-time approximation, collisions relax the density matrix toward a reference state at the rate $\gamma_{\mathrm{rel}}$:
\begin{equation*}
  \ii\,\partial_t\,\delta\varrho_{\vect{k}} = -D_{\vect{k}}\,\delta\varrho_{\vect{k}} + (n_{\vect{k}} - n_{\vect{k}+\vect{q}})\,\phi_{\mathrm{ext}} - \ii\gamma_{\mathrm{rel}}\bigl(\delta\varrho_{\vect{k}} - \delta\varrho^{\mathrm{loc}}_{\vect{k}}\bigr) .
\end{equation*}
With the global equilibrium as the reference, $\delta\varrho^{\mathrm{loc}}_{\vect{k}} = 0$, the solution is the collisionless solution with $\omega$ replaced by $\omega + \ii\gamma_{\mathrm{rel}}$, the naive substitution. Mermin's reference is the equilibrium state in which the density already has its instantaneous value, since collisions conserve the number of electrons (\cref{sec:4.5}). A local equilibrium with a local chemical potential $\mu_{\mathrm{ch}} + \delta\mu(\vect{r}, t)$ is the equilibrium state in the static potential $-\delta\mu$, and we write this potential as $\phi_{\mathrm{loc}}(t)\,\ee^{\ii\vect{q}\cdot\vect{r}} + \text{c.c.}$ The density matrix of the local equilibrium is the time-independent solution of the collisionless equation in a static potential,
\begin{equation*}
  \delta\varrho^{\mathrm{loc}}_{\vect{k}} = \frac{n_{\vect{k}} - n_{\vect{k}+\vect{q}}}{D_{\vect{k}}}\,\phi_{\mathrm{loc}}(t), \qquad
  \frac{1}{\Vol}\sum_{\vect{k}}\delta\varrho^{\mathrm{loc}}_{\vect{k}} = R^0(\vect{q},0)\,\phi_{\mathrm{loc}}(t) ,
\end{equation*}
since $(1/\Vol)\sum_{\vect{k}}(n_{\vect{k}} - n_{\vect{k}+\vect{q}})/D_{\vect{k}}$ is the static bubble. The static response that enters here is the isothermal one, the value of the Matsubara response at $\omega^{\mathrm{B}}_0 = 0$ in \cref{sec:4.7}, which for free electrons at $\vect{q} \ne 0$ coincides with $R^0(\vect{q},\omega \to 0)$. The amplitude $\phi_{\mathrm{loc}}$ is fixed by requiring the local-equilibrium state to have the actual density,
\begin{equation}
  \frac{1}{\Vol}\sum_{\vect{k}}\delta\varrho^{\mathrm{loc}}_{\vect{k}} = \frac{1}{\Vol}\sum_{\vect{k}}\delta\varrho_{\vect{k}} , \qquad\text{that is}\qquad
  \phi_{\mathrm{loc}} = \frac{\delta\nd}{R^0(\vect{q},0)} .
  \label{eq:D.9-condition}
\end{equation}

For a field at frequency $\omega$, the equation of motion becomes algebraic,
\begin{equation*}
  (\omega + \ii\gamma_{\mathrm{rel}} + D_{\vect{k}})\,\delta\varrho_{\vect{k}} = (n_{\vect{k}} - n_{\vect{k}+\vect{q}})\Bigl[\phi_{\mathrm{ext}} + \frac{\ii\gamma_{\mathrm{rel}}\,\phi_{\mathrm{loc}}}{D_{\vect{k}}}\Bigr] .
\end{equation*}
We divide by the factor on the left and use the partial fractions $1/[D(\omega + \ii\gamma_{\mathrm{rel}} + D)] = (\omega + \ii\gamma_{\mathrm{rel}})^{-1}[1/D - 1/(\omega + \ii\gamma_{\mathrm{rel}} + D)]$; the sum over $\vect{k}$ then produces two bubbles, one at the complex frequency $\omega + \ii\gamma_{\mathrm{rel}}$ and one static. With the shorthand $R^0_\gamma = R^0(\vect{q},\omega + \ii\gamma_{\mathrm{rel}})$ and $R^0_{\mathrm{st}} = R^0(\vect{q},0)$,
\begin{equation*}
  \delta\nd = R^0_\gamma\,\phi_{\mathrm{ext}} + \frac{\ii\gamma_{\mathrm{rel}}}{\omega + \ii\gamma_{\mathrm{rel}}}\bigl(R^0_{\mathrm{st}} - R^0_\gamma\bigr)\,\phi_{\mathrm{loc}} .
\end{equation*}
Inserting \cref{eq:D.9-condition} and collecting the terms in $\delta\nd$ on the left gives
\begin{equation*}
  \frac{\omega + \ii\gamma_{\mathrm{rel}}R^0_\gamma/R^0_{\mathrm{st}}}{\omega + \ii\gamma_{\mathrm{rel}}}\,\delta\nd = R^0_\gamma\,\phi_{\mathrm{ext}} ,
\end{equation*}
and the response $\delta\nd/\phi_{\mathrm{ext}}$ is
\begin{equation}
  R^{\mathrm{Mer}}(\vect{q},\omega) = \frac{(1 + \ii\gamma_{\mathrm{rel}}/\omega)\,R^0(\vect{q},\omega + \ii\gamma_{\mathrm{rel}})}{1 + (\ii\gamma_{\mathrm{rel}}/\omega)\,R^0(\vect{q},\omega + \ii\gamma_{\mathrm{rel}})/R^0(\vect{q},0)} .
  \label{eq:D.9-R}
\end{equation}
In the random-phase approximation the Mermin dielectric function is $\varepsilon_{\mathrm{Mer}} = 1 - v_{\mathrm{C}}R^{\mathrm{Mer}}$. With $\varepsilon_{\mathrm{Lin}} = 1 - v_{\mathrm{C}}R^0$ for the Lindhard dielectric function, so that $R^0_\gamma/R^0_{\mathrm{st}} = [\varepsilon_{\mathrm{Lin}}(\vect{q},\omega + \ii\gamma_{\mathrm{rel}}) - 1]/[\varepsilon_{\mathrm{Lin}}(\vect{q},0) - 1]$, the Mermin dielectric function takes the form given by Mermin~\cite{mermin1970lindhard},
\begin{equation}
  \varepsilon_{\mathrm{Mer}}(\vect{q},\omega) = 1 + \frac{(1 + \ii\gamma_{\mathrm{rel}}/\omega)\,\bigl[\varepsilon_{\mathrm{Lin}}(\vect{q},\omega + \ii\gamma_{\mathrm{rel}}) - 1\bigr]}{1 + (\ii\gamma_{\mathrm{rel}}/\omega)\,\bigl[\varepsilon_{\mathrm{Lin}}(\vect{q},\omega + \ii\gamma_{\mathrm{rel}}) - 1\bigr]/\bigl[\varepsilon_{\mathrm{Lin}}(\vect{q},0) - 1\bigr]} .
  \label{eq:D.9-eps}
\end{equation}
Nothing in the derivation used $T = 0$; at finite temperature $R^0$ and $\varepsilon_{\mathrm{Lin}}$ are taken at that temperature.

\subsection*{Causality}

The bubble at the shifted frequency, $R^0(\vect{q},z + \ii\gamma_{\mathrm{rel}})$, is analytic for $\Imag z > -\gamma_{\mathrm{rel}}$. The analyticity of the shifted bubble alone does not make $R^{\mathrm{Mer}}$ causal, since the denominator of \cref{eq:D.9-R}, which is proportional to $z + \ii\gamma_{\mathrm{rel}}R^0_\gamma/R^0_{\mathrm{st}}$, could vanish in the upper half plane, and a pole there would be a response that grows after the field is gone. The relaxation equation has no growing mode, and the proof uses only its structure.

A zero of the denominator is a free oscillation. With $\phi_{\mathrm{ext}} = 0$ and the time dependence $\ee^{-\ii zt}$, the equation for $\delta\nd$ obtained above by collecting terms reads $(z + \ii\gamma_{\mathrm{rel}}R^0_\gamma/R^0_{\mathrm{st}})\,\delta\nd = 0$, so a nonzero density exists exactly where the denominator vanishes. The nonzero density comes with the nonzero coherences $\delta\varrho_{\vect{k}} = (n_{\vect{k}} - n_{\vect{k}+\vect{q}})\,\ii\gamma_{\mathrm{rel}}\phi_{\mathrm{loc}}/[D_{\vect{k}}(z + \ii\gamma_{\mathrm{rel}} + D_{\vect{k}})]$. A zero with $\Imag z > 0$ would be a field-free solution of the equation of motion that grows as $\ee^{(\Imag z)t}$.

No growing field-free solution exists. Write $w_{\vect{k}} = (n_{\vect{k}} - n_{\vect{k}+\vect{q}})/(\xi_{\vect{k}+\vect{q}} - \xi_{\vect{k}})$, which is positive wherever the two occupations differ, because $n_{\mathrm{F}}$ falls with energy, and zero where they are equal. Let $\varsigma = \sum_{\vect{k}}\delta\varrho_{\vect{k}}$ and $w_{\mathrm{tot}} = \sum_{\vect{k}}w_{\vect{k}}$. Since $R^0_{\mathrm{st}} = -w_{\mathrm{tot}}/\Vol$, the condition \cref{eq:D.9-condition} makes the reference state $\delta\varrho^{\mathrm{loc}}_{\vect{k}} = w_{\vect{k}}\,\varsigma/w_{\mathrm{tot}}$, and the equation of motion without a field is
\begin{equation*}
  \partial_t\,\delta\varrho_{\vect{k}} = \ii D_{\vect{k}}\,\delta\varrho_{\vect{k}} - \gamma_{\mathrm{rel}}\Bigl(\delta\varrho_{\vect{k}} - w_{\vect{k}}\frac{\varsigma}{w_{\mathrm{tot}}}\Bigr) .
\end{equation*}
A coherence with $w_{\vect{k}} = 0$ obeys $\partial_t\,\delta\varrho_{\vect{k}} = (\ii D_{\vect{k}} - \gamma_{\mathrm{rel}})\,\delta\varrho_{\vect{k}}$ and decays at the rate $\gamma_{\mathrm{rel}}$, so a solution $\propto \ee^{-\ii zt}$ with $\Imag z > 0$ has every coherence with $w_{\vect{k}} = 0$ equal to zero, and $\varsigma$ is a sum over the coherences with positive $w_{\vect{k}}$. For these coherences take the weighted norm
\begin{equation*}
  \Phi_\varrho(t) = \sum_{w_{\vect{k}} > 0}\frac{|\delta\varrho_{\vect{k}}|^2}{w_{\vect{k}}} \ge 0 .
\end{equation*}
The term $\ii D_{\vect{k}}$ rotates each phase and does not change $\Phi_\varrho$, and
\begin{equation*}
  \frac{\dd\Phi_\varrho}{\dd t} = -2\gamma_{\mathrm{rel}}\Bigl[\Phi_\varrho - \frac{|\varsigma|^2}{w_{\mathrm{tot}}}\Bigr] \le 0 ,
\end{equation*}
because the Cauchy--Schwarz inequality gives $|\varsigma|^2 = |\sum_{\vect{k}}(\delta\varrho_{\vect{k}}/w_{\vect{k}}^{1/2})\,w_{\vect{k}}^{1/2}|^2 \le \Phi_\varrho w_{\mathrm{tot}}$. A growing solution would make $\Phi_\varrho \propto \ee^{2(\Imag z)t}$ grow, so none exists, and the denominator of \cref{eq:D.9-R} has no zero with $\Imag z > 0$. At $z = 0$ the numerator and the denominator both contain $\ii\gamma_{\mathrm{rel}}/z$, the singularity is removable, and $R^{\mathrm{Mer}} \to R^0(\vect{q},0)$. With the decay at large $|z|$ shown below, $R^{\mathrm{Mer}}$ is analytic in the upper half plane and falls off at infinity, a causal response in the sense of \cref{app:C.5}.

\subsection*{Particle conservation}

The sum of the equation of motion over $\vect{k}$ is the continuity equation, and the reference state decides whether it holds. In the term with the potential, $\sum_{\vect{k}}(n_{\vect{k}} - n_{\vect{k}+\vect{q}}) = 0$, because the two sums run over the same occupations. In the term with $D_{\vect{k}}$, the free dispersion gives $-D_{\vect{k}} = \xi_{\vect{k}+\vect{q}} - \xi_{\vect{k}} = \vect{q}\cdot(\vect{k} + \vect{q}/2)/m_e$. The factor $(\vect{k} + \vect{q}/2)/m_e$ is the plane-wave matrix element $\langle\vect{k} + \vect{q}|\tfrac12\{\vect{p}/m_e, \ee^{\ii\vect{q}\cdot\vect{r}}\}|\vect{k}\rangle$ of the paramagnetic current $\vect{j}_e(\vect{q}) = \sum_i\tfrac12\{\vect{p}_i/m_e, \ee^{\ii\vect{q}\cdot\vect{r}_i}\}$, the velocity of the electrons weighted by the plane wave (\cref{sec:4.3}). By \cref{eq:2.3-onebody}, $\vect{j}_e(\vect{q})^\dagger = \sum_{\vect{k}}[(\vect{k} + \vect{q}/2)/m_e]\,c^\dagger_{\vect{k}}c_{\vect{k}+\vect{q}}$, so that $-\sum_{\vect{k}}D_{\vect{k}}\delta\varrho_{\vect{k}} = \vect{q}\cdot\delta\langle\vect{j}_e(\vect{q})^\dagger\rangle$. At frequency $\omega$ the summed equation reads
\begin{equation}
  \omega\,\delta\langle n_e(\vect{q})^\dagger\rangle - \vect{q}\cdot\delta\langle\vect{j}_e(\vect{q})^\dagger\rangle = -\ii\gamma_{\mathrm{rel}}\sum_{\vect{k}}\bigl(\delta\varrho_{\vect{k}} - \delta\varrho^{\mathrm{loc}}_{\vect{k}}\bigr) .
  \label{eq:D.9-continuity}
\end{equation}
Without collisions the right-hand side is zero, and \cref{eq:D.9-continuity} is the Fourier transform of $\partial_t\nd_e + \nabla\cdot\vect{j} = 0$: the density at a crest changes only as fast as the current carries electrons in or out. The naive substitution, with $\delta\varrho^{\mathrm{loc}} = 0$, leaves $-\ii\gamma_{\mathrm{rel}}\,\delta\langle n_e(\vect{q})^\dagger\rangle$ on the right, a density that decays at the rate $\gamma_{\mathrm{rel}}$ whatever the current does. $R^{\mathrm{Mer}}$ conserves particle number in the sense that Mermin's condition \cref{eq:D.9-condition} sets the right-hand side to zero exactly, at every frequency and momentum. Collisions now change only the current, and through the current the density.

\subsection*{Long-wavelength limit}

For $q \to 0$ at fixed $\omega$ and $\gamma_{\mathrm{rel}}$, the frequency $\omega + \ii\gamma_{\mathrm{rel}}$ lies far outside the continuum, $|\omega + \ii\gamma_{\mathrm{rel}}| \gg qv_{\mathrm{F}}$, and the high-frequency form \cref{eq:4.5-moments} gives $R^0_\gamma \to \nd_eq^2/[m_e(\omega + \ii\gamma_{\mathrm{rel}})^2]$, while $R^0_{\mathrm{st}}$ stays finite. The ratio $R^0_\gamma/R^0_{\mathrm{st}}$ vanishes as $q^2$, the denominator of \cref{eq:D.9-R} tends to one, and $R^{\mathrm{Mer}} \to \nd_eq^2/[m_e\,\omega(\omega + \ii\gamma_{\mathrm{rel}})]$, which with $\varepsilon_{\mathrm{Mer}} = 1 - v_{\mathrm{C}}R^{\mathrm{Mer}}$ is the Drude function of \cref{sec:4.5}. At large $\omega$ the same form gives $R^{\mathrm{Mer}} \to \nd_eq^2/m_e\omega^2$, the leading tail of the naive substitution as well, and by \cref{eq:C.5-highfreq} the Mermin response and the naive substitution both obey $\int_0^\infty\dd\omega\,\omega\,\Imag R = -\pi\nd_eq^2/2m_e$.
